% !TeX root = ../深度学习与大语言模型技术原理.tex
% 由主文档通过 \input 载入；请编译主文档。
\setcounter{chapter}{23} % 保留原书章号。
\chapter{量化、缓存压缩与剪枝}
\label{ch:quantization}

\section{量化映射与误差来源}

量化（quantization）用有限码字近似实数。均匀仿射量化可写为
\begin{equation}
q=\operatorname{clip}\!\left(\operatorname{round}(x/s)+z,q_{\min},q_{\max}\right),\qquad
\widehat x=s(q-z),\quad s>0,
\label{eq:quantization}
\end{equation}
其中 $s$ 为尺度，$z$ 为零点。未发生裁剪且采用最近舍入时，$|x-\widehat x|\le s/2$；发生裁剪时还包含饱和误差。逐张量、逐通道和分组尺度在元数据成本与局部适配之间折中。

权重量化、激活量化与 KV 量化作用于不同对象。W4A16 常表示4位权重、16位激活，但不完整说明累积精度、尺度、零点或内核实现。训练后量化（post-training quantization, PTQ）使用已训练模型和必要校准数据；量化感知训练（quantization-aware training, QAT）在训练中模拟量化误差，常用直通估计器近似处理不可微舍入，其梯度并非舍入函数的真实导数。

低位存储降低字节数，实测加速仍取决于硬件是否支持相应运算、反量化成本、矩阵形状和内存布局。某些权重量化主要节约带宽；若任务已计算受限或小算子占主导，位数减半不意味着延迟减半。

\section{GPTQ 的重构目标与 AWQ 的缩放}

令校准输入 $X\in\mathbb R^{d_i\times N}$，GPTQ 类方法关注层输出重构
\[
\min_{\widehat W\in\mathcal Q}\|(W-\widehat W)X\|_F^2,
\]
其中 $\mathcal Q$ 是量化约束集合。每个输出行的二次型包含 $XX^\top$，因此可用二阶信息估计量化某一权重后对其余权重的补偿。GPTQ 通过逐步量化与误差补偿近似求解，通常加阻尼改善病态矩阵的数值稳定性\cite{gptq}。这是校准分布上的局部重构，不是端到端任务损失的全局最优解。

激活感知权重量化（Activation-aware Weight Quantization, AWQ）依据激活统计识别重要通道，并搜索有利的缩放\cite{awq}。对可逆对角矩阵 $S$，有 $Wx=(WS)(S^{-1}x)$；精确算术中函数不变，但对 $WS$ 量化后，误差作用变为 $(\widehat{WS}-WS)S^{-1}x$，从而改变各通道的量化难度。工程上要把等价缩放合理吸收到相邻算子，避免额外开销。

两类方法都依赖校准样本代表性。短通用语料上的低重构误差，不保证长上下文、代码或特殊工具格式不退化。至少应比较困惑度、目标任务质量、长上下文稳定性和服务延迟。

\section{FP8、FP4 与缓存量化}

FP8、FP4 指低位浮点格式家族，具体指数、尾数、尺度和特殊值约定可能不同。整数格式在区间内均匀分配刻度，低位浮点格式的间隔随数值范围变化；两者都需要考虑离群值。块尺度可以提高局部动态范围适配，但增加尺度数据和运算。TensorRT-LLM 等系统的支持取决于所选后端与硬件，应按精度文档核对\cite{trt}。

KV Cache 量化直接减少式~\eqref{eq:kv-memory} 的元素字节数。键的误差改变查询--键匹配，从而改变 softmax 权重；值的误差改变被加权内容，两者传播路径不同。逐通道或逐词元尺度、分组大小、近期高精度残留区等策略有不同代价，不能仅凭权重量化方案推断缓存量化效果。

长上下文会让缓存误差影响更多后续查询。评测应包括长序列、不同位置证据、不同批量和持续生成，而非只测短输入的首个词元。尺度元数据、反量化临时空间与带宽也必须计入实际容量和速度。

\section{结构化与非结构化剪枝}

剪枝（pruning）删除权重或计算结构。非结构化剪枝可保留任意稀疏模式，幅值剪枝按 $|w|$ 选择小权重；若输入方差不同，幅值小不一定输出影响小，式~\eqref{eq:activation-error} 说明了原因。稀疏矩阵还需存索引，普通稠密内核不会因大量零元素自动加速。

结构化剪枝删除通道、注意力头、FFN 中间维度或整层，能够改变实际矩阵形状，更容易被常规内核利用，但通常更强地改变表示能力。半结构化稀疏规定固定小组内保留数量，需要匹配硬件支持。剪枝后常需恢复训练或蒸馏，并重新检查残差维度和归一化参数。

量化减少每个数的表示成本，剪枝减少部分参数或运算，蒸馏训练新的学生行为；三者可以组合，但误差不一定简单相加。选择顺序和恢复策略应通过质量--延迟--内存的联合比较确定。

